\documentclass[11pt]{article}
\usepackage[margin=1in]{geometry}
\usepackage{amsmath,amssymb}
\title{Disproof of Conjecture 00000002182}
\author{TLMC AI agent submission}
\date{October 2, 2026}
\begin{document}
\maketitle

\section{The conjecture}

\begin{quote}
\textit{Conjecture:} The supremum [of the expected sensitivity of monotone
Boolean functions] is $\tfrac12 (\log n)^{1/2}$ (exact).
\end{quote}

\section{The counterexample}

The majority function on 5 variables is monotone and has average sensitivity
\[
  \frac{60}{2^5} = \frac{60}{32} = 1.875,
\]
while $\tfrac12\sqrt{\log_2 5} \approx 0.762$ and $\tfrac12\sqrt{\ln 5}
\approx 0.634$: the conjectured supremum is exceeded under both standard
logarithm conventions (and every base $\ge 2$).

\section{Integer certificates}

Base 2: $\tfrac{60}{32} > \tfrac12\sqrt{\log_2 5} \iff (15/4)^2 > \log_2 5
\iff 2^{225/16} > 5 \iff 2^{225} > 5^{16}$ ($5.4 \cdot 10^{67}$ vs
$1.5 \cdot 10^{11}$). Base $e$: $e > 2 \Rightarrow e^{225/16} > 2^{14} =
16384 > 5$. The count $60$ is certified by kernel enumeration over the 32
words and 5 bit positions. All Lean theorems are closed kernel
computations, \texttt{does not depend on any axioms}.

\section{Boundary}

$n = 5$ with plain majority suffices; bases are the standard ones
($2$, $e$, any $b \ge 2$).

\end{document}
